% SHARD-ID: CA-03-INDEFINITE-PARTICLE-FOCK
% SHARD-TITLE: Indefinite Particle Number: the Full Fock Completion
% SHARD-SUMMARY: Makes the indefinite-particle Hilbert space rigorous as the l2 (Hilbert) direct sum of completed tensor powers, with the vacuum at degree zero.
% SHARD-SUMMARY: Distinguishes the algebraic direct sum (finite particle number, dense) from its Hilbert completion (indefinite particle number) and shows the Fock space is a Hilbert space with mutually orthogonal sectors.
% SHARD-SUMMARY: Identifies this as the full / distinguishable Fock space and locates Bose/Fermi statistics as the deferred per-sector permutation quotient.
% SHARD-KEYWORDS: Fock space, indefinite particle number, Hilbert direct sum, completed tensor product, vacuum, distinguishable Fock, second quantization

\section{Indefinite Particle Number: the Full Fock Completion}
\label{sec:indefinite-particle-fock}

\subsection{Sources and Dependencies}

This shard completes CA-02's finite AND/OR/vacuum grammar to a genuine
\emph{indefinite}-particle Hilbert space. It depends on CA-02 (the two monoidal
structures, the units $\zero$ and $\mathbb C$, the carrier $\mathcal V$ and grading
$\mathcal V_n$). The named Fock construction is anchored locally:

% Source: references/cft/Schottenloher2008/Schottenloher2008.md:4216
%   "D:=\bigoplus_{N=0}^{\infty} H_N"
% Source: references/cft/Schottenloher2008/Schottenloher2008.md:4218
%   "(H_0 = C with the vacuum Omega := 1 in H_0)"
% Source: references/cft/Schottenloher2008/Schottenloher2008.md:4222
%   "The completion of D with respect to this inner product is the Fock space H."
% Source: literature/md/1910.10619/1910.10619.md:57
%   "F(C^d) = \bigoplus_{N=0}^{\infty} H_N = \bigoplus_{N=0}^{\infty}(C^d \otimes ... \otimes C^d)  (N factors)"
% Source: literature/md/1910.10619/1910.10619.md:59
%   "the *vacuum*, is assigned the Hilbert space (C^d)^{\otimes 0} = C."
% Source: literature/md/2010.11121/2010.11121.md:258
%   "bosonic Fock space over h: F_+(h) = C\Omega \oplus \bigoplus_n h^{\otimes_s n}".

The sibling project's finite grammar (\texttt{AQM-70}) is recovered as the
finite-truncation of the object built here.

\subsection{Two Direct Sums --- the Load-Bearing Distinction}

Let $\{W_n\}_{n\ge 0}$ be a countable family of Hilbert spaces.
\begin{itemize}
  \item The \textbf{algebraic direct sum} $\bigoplus^{\mathrm{alg}}_{n}W_n$ consists
  of finite-support sequences $(w_n)$ ($w_n\in W_n$, $w_n=0$ for all but finitely
  many $n$). It is a pre-Hilbert space, generally \emph{incomplete}.
  \item The \textbf{Hilbert ($\ell^2$) direct sum} $\hoplus_{n}W_n$ consists of
  sequences with $\sum_n\|w_n\|_{W_n}^2<\infty$, with inner product
  $\langle (w_n),(v_n)\rangle=\sum_n\langle w_n,v_n\rangle_{W_n}$ (absolutely
  convergent by Cauchy--Schwarz). It is \emph{complete}, hence a Hilbert space.
\end{itemize}
$\bigoplus^{\mathrm{alg}}_{n}W_n$ is a \emph{dense} subspace of $\hoplus_{n}W_n$,
and the latter is its completion (standard functional analysis).

\paragraph{Why indefinite particle number forces $\hoplus$.}
A state with nonzero amplitude in \emph{unboundedly} many particle-number sectors
(a genuine superposition over indefinitely many particles) has infinite support: it
lies in $\hoplus$ but not in $\bigoplus^{\mathrm{alg}}$. ``Indefinite particle
number'' $\Leftrightarrow$ ``not confined to finitely many sectors''
$\Rightarrow$ the $\ell^2$ completion is \emph{required}. This is exactly what
CA-02's finite grammar lacks.

\subsection{Completed Tensor Powers}

For Hilbert spaces $H,K$, the \emph{completed} (Hilbert) tensor product
$H\hotimes K$ is the completion of the algebraic tensor product under the inner
product determined by $\langle a\otimes b,\,c\otimes d\rangle
=\langle a,c\rangle_H\langle b,d\rangle_K$ (well defined and positive definite).
For finite-dimensional factors $\hotimes=\otimes$; the completion bites only for
infinite-dimensional one-particle spaces. Define the \emph{completed tensor powers}
\[
  H^{\hotimes 0} := \mathbb C\ \ (\text{the vacuum sector}),
  \qquad
  H^{\hotimes n} := \underbrace{H\hotimes\cdots\hotimes H}_{n},
\]
the bracketing immaterial up to the canonical unitary associativity isomorphism.

\subsection{The Full Fock Space}

\paragraph{Definition.}
For a one-particle Hilbert space $H$, the \emph{full Fock space} is
\[
  \boxed{\ \Fock(H)\;:=\;\hoplus_{n\ge 0} H^{\hotimes n}\ },
  \qquad H^{\hotimes 0}=\mathbb C,
\]
i.e.\ $\psi=(\psi_n)_{n\ge0}$ with $\psi_n\in H^{\hotimes n}$ and
$\|\psi\|^2=\sum_n\|\psi_n\|^2<\infty$. The unit vector $\vac=(1,0,0,\dots)$ in the
degree-$0$ sector $\mathbb C$ is the \emph{vacuum}. This is the
\emph{distinguishable} Fock space (the author's term in the local source
\texttt{1910.10619}); in operator algebra it is also called the \emph{free} or
\emph{Boltzmann} Fock space (standard terminology, not anchored in a local
source), namely the Hilbert completion of the tensor algebra
$T(H)=\bigoplus_n H^{\otimes n}$. It is \emph{not} the symmetric (bosonic)
or antisymmetric (fermionic) Fock space; see ``Exchange sectors come later''
below.

\paragraph{Compiler status.}
In the compiler grammar this is not an additional primitive constructor.  It is
the Hilbert completion of the ordinary tensor-power expressions
\(\one,H,H^{\hotimes2},\ldots\): a countable OR over all finite AND powers.  If a
one-particle symmetry representation is present, CA-05 makes this same object a
graded direct-sum representation.

\paragraph{Theorem (Fock is a Hilbert space).}
$\Fock(H)$ is a Hilbert space; (i) the algebraic direct sum
$\bigoplus^{\mathrm{alg}}_n H^{\hotimes n}$ (the finite-particle vectors) is
\emph{dense}; (ii) distinct-$n$ sectors are \emph{mutually orthogonal},
$\langle\psi,\varphi\rangle=\sum_n\langle\psi_n,\varphi_n\rangle$, so
$H^{\hotimes m}\perp H^{\hotimes n}$ for $m\neq n$; (iii) on the degree-$0$ sector
the inner product is the standard one on $\mathbb C$ and $\vac$ is a unit vector.
\emph{Proof.} An instance of the $\ell^2$-direct-sum facts above with
$W_n=H^{\hotimes n}$: completeness, density of the algebraic sum, and summand
orthogonality are exactly those statements. $\square$

\paragraph{The $1/N!$ caveat (a convention trap).}
Schottenloher's inner product on $D=\bigoplus_N H_N$ weights sector $N$ by $1/N!$;
that factor is a \emph{symmetric}-Fock normalisation correcting the $N!$
overcounting of symmetric tensors. The full Fock space above uses the \emph{plain}
$\ell^2$ sum with \emph{no} combinatorial weight. Do not copy the $1/N!$ into the
distinguishable case.
% Source: references/cft/Schottenloher2008/Schottenloher2008.md:4220
%   inner product "f_0 g_0 + sum_{N>=1} (1/N!) <f_N,g_N>" -- symmetric-Fock weight.

\subsection{The Grammar Generates the Indefinite-Particle Space}

Let the one-particle menu over a countable type set $X$ be
$P=\mathrm{OR}_{x\in X}K_x$, with carrier $\mathcal V(P)=\hoplus_{x\in X}K_x$ (the
$\ell^2$ direct sum; an ordinary finite direct sum when $X$ is finite). The
indefinite-particle space over this menu is $\Fock(\mathcal V(P))$. Expanding each completed tensor power by distributivity
(CA-02's $\delta$, extended unitarily to the completions) gives the
$(n,\text{type-word})$-graded decomposition
\[
  \mathcal V(P)^{\hotimes n}\ \cong\ \hoplus_{w\in X^n} K_w,
  \quad K_w:=K_{x_1}\hotimes\cdots\hotimes K_{x_n},
  \qquad\Longrightarrow\qquad
  \Fock(\mathcal V(P))\ \cong\ \hoplus_{n\ge 0}\ \hoplus_{w\in X^n} K_w .
\]
Each of the $n$ tensor slots independently chooses a type $x_i\in X$ (two slots may
choose the same type; nothing is symmetrised). The grading is $\mathcal V_n=\hoplus_{w\in X^n}K_w$ ($n$ = particle number, $w\in X^n$ = ordered type-word) ---
the $\ell^2$-completed, all-$n$ analogue of CA-02's finite grade-$n$ sector. The
sibling project's finite expressions ($P_X^{\,\mathrm{AND}\,m}$, the truncation
$T_{\le M}$) are precisely the finite-support / finite-$M$ restrictions of this
object; that is the exact harmonisation: \texttt{AQM-70} is the
$\bigoplus^{\mathrm{alg}}$/finite-$M$ truncation of $\Fock$.

\subsection{Exchange Sectors Come Later}

The full Fock space already carries symmetry representations when \(H\) does:
internal symmetries act sectorwise by tensor powers (CA-05), and for each \(n\)
the permutation group \(S_n\) acts unitarily on \(H^{\hotimes n}\) by permuting
tensor slots.  What is deferred is the \emph{sector selection}.  Bosonic and
fermionic Fock spaces are the images of the per-sector orthogonal projectors
\[
  \Gamma_{\mathrm s}(H)=P_{\mathrm{sym}}\,\Fock(H)
  \quad(\text{symmetric / Bose}),
  \qquad
  \Gamma_{\mathrm a}(H)=P_{\mathrm{anti}}\,\Fock(H)
  \quad(\text{antisymmetric / Fermi}),
\]
with para-statistics the other Young-symmetry isotypic components. This is the
``impose an equivalence relation under particle exchange'' step.
% Source: literature/md/1910.10619/1910.10619.md:41
%   "one can obtain Bose-/Fermi-Fock space by imposing an equivalence relation
%    under particle exchange."
We define no projector and prove nothing here: the construction of these symmetry
quotients/subspaces, and the genuinely anyonic case (where standard Fock space is
not the right object), are deferred to later shards.

\subsection{Lamport Dependency Skeleton}

{\small
\[
\setlength{\arraycolsep}{2pt}
\begin{array}{rcl}
D6 &:& \bigoplus^{\mathrm{alg}}\!(\text{finite support})\ \subset\
       \hoplus\,(\ell^2):\ \langle(w_n),(v_n)\rangle=\textstyle\sum_n\langle w_n,v_n\rangle.\\
D7 &:& \text{completed tensor power } H^{\hotimes n},\ H^{\hotimes 0}=\mathbb C\ (\text{vacuum}).\\
D8 &:& \Fock(H):=\hoplus_{n\ge0}H^{\hotimes n}\ \text{(full / distinguishable Fock);}\ \vac=(1,0,\dots).\\
D9 &:& P=\mathrm{OR}_{x\in X}K_x;\quad \Fock(\mathcal V(P))=\hoplus_{n\ge0}\mathcal V(P)^{\hotimes n}.\\
T3 &:& \text{(D6,D7) } \Fock(H)\ \text{Hilbert; } \bigoplus^{\mathrm{alg}}\ \text{dense; }
       H^{\hotimes m}\perp H^{\hotimes n}\ (m\neq n).\\
T4 &:& \text{indefinite particle no.} \Leftrightarrow \text{infinite support} \Rightarrow \hoplus\ \text{(not } \bigoplus^{\mathrm{alg}}).\\
T5 &:& \text{(D8,}\delta\text{) } \Fock(\mathcal V(P))\cong\hoplus_{n\ge0}\hoplus_{w\in X^n}K_w\ ((n,\text{type-word})\text{ grading}).\\
T6 &:& \Gamma_{\mathrm s}=P_{\mathrm{sym}}\Fock,\ \Gamma_{\mathrm a}=P_{\mathrm{anti}}\Fock\ \text{deferred (statistics later)}.\\
\end{array}
\]
}

\subsection{Scope and Pitfalls}

\textbf{In scope:} the indefinite-particle direction, i.e.\ countable
$\mathrm{OR}$ (the Fock completion). \textbf{Out of scope:} countable
$\mathrm{AND}$ (an \emph{infinite tensor product}), which is analytically a
different object (von Neumann's infinite tensor product, needing a stabilising
reference sequence) --- not developed here.

Pitfalls: (i) \emph{algebraic} $\bigoplus^{\mathrm{alg}}$ (finite particle number,
dense, incomplete) vs.\ \emph{completed} $\hoplus$ (indefinite, complete) --- the
distinction this shard exists to make. (ii) For infinite-dimensional one-particle
$H$, the $\mathrm{AND}$ of CA-02 must mean the \emph{completed} $\hotimes$, not the
algebraic $\otimes$. (iii) The $1/N!$ symmetric-Fock weight must not be imported
into the distinguishable case. (iv) All sector identities hold on dense subspaces
and extend by continuity; states are $\ell^2$-summable and cross-sector inner
products vanish.
